\scanpage{21}
% archive sheet number: 20
\unsure{Donc on a|D'où} les relations
\[
\left\{\begin{aligned}
&\widetilde\sigma_0^2=\widetilde\sigma_1^2=\widetilde\sigma_\infty^2=1\\
&\widetilde\sigma_0\widetilde\sigma_1=\widetilde\sigma_1\widetilde\sigma_\infty=\widetilde\sigma_\infty\widetilde\sigma_0=\deleted{\widetilde\rho}\rho\deleted{(\unsure{=\rho})}\\
&(\deleted{\widetilde\rho}\rho^3=1\text{ est conséquence des précédentes})
\end{aligned}\right.\tag*{(13)}
\]
% slanted
\margin{NB $\rho\widetilde\sigma_i\rho^{-1}=\widetilde\sigma_{\unsure{\dot\rho i|\dot\rho(i)}}$ i.e. $\rho\widetilde\sigma_0\rho^{-1}=\widetilde\sigma_1$, $\rho\widetilde\sigma_1\rho^{-1}=\widetilde\sigma_\infty$, $\rho\widetilde\sigma_\infty\rho^{-1}=\widetilde\sigma_0$.}
et l'image réciproque $\unsure{\widetilde{\mathscr E}|\widetilde{\mathcal E}}$ de $\varphi(\unsure{\widetilde{\mathfrak S}_3|\mathfrak S_3|\widetilde{\mathscr S}_3|\widetilde G_3})\subset\Gamma$ (s-groupe d'indice 2) est produit $\frac12$ direct de $\widetilde{\mathfrak S}_3$ et de $\pi$. Il faut expliciter les opérations des $\widetilde\sigma_a$ sur $\pi$ pour avoir la structure complète du groupe. On trouve que $\widetilde\sigma_a$ est au-dessus de $\varphi(\deleted{\unsure{\widetilde\sigma_a}}\dot\sigma_a)$ i.e. la symétrie \unsure{p.r.|p.v.} au plan passant par $P,\overline P,\deleted{\unsure{a}}\unsure{i}$, qui transforme $l_0$ en $l_0^{-1}$, $\deleted{l}l_\infty$ en $l_1^{-1}$, $l_1$ en $l_\infty^{-1}$ :
% slanted
\margin{\unsure{Ce sont elles obtenues par transport de structure}, via les opérations des \unsure{mêmes} (4,8).}
\[
\left\{\begin{aligned}
\widetilde\sigma_0l_0\widetilde\sigma_0^{-1}&=l_0^{-1}\qquad
\widetilde\sigma_0l_1\widetilde\sigma_0^{-1}=l_\infty^{-1}\qquad
\widetilde\sigma_\infty l_\infty\widetilde\sigma_0^{-1}=l_1^{-1}\\
\widetilde\sigma_1l_1\widetilde\sigma_1^{-1}&=l_1^{-1}\qquad\cdots\\
&\cdots
\end{aligned}\right.\tag*{(14)}
\]
% slanted
\margin{\uline{Critique}. La notation $\widetilde\sigma_0$ suggère à tort que $\widetilde\sigma_0$ est au-dessus de $\sigma_0$, il vaut mieux écrire $\dot\sigma_i'=\dot\sigma_i\tau$ et noter $\widetilde\sigma_i'$ au lieu des $\widetilde\sigma_i$ $(i\in\{0,1,\infty\})$, ou mieux $\sigma_i'$ (sans $\sim$) --- pour réserver $\widetilde\sigma_i$ à des éléments de $\unsure{S\mathscr E|S_\Sigma|S\Sigma}$, cf plus bas\unsure{...}}

Les relations (13), (14), plus la relation fondamentale $l_\infty l_1l_0=1$ (1), font une présentation de $\unsure{\widetilde{\mathscr G}|\widetilde{\mathscr E}|\widetilde G}$, dont on peut si on veut éliminer $\rho$, (en omettant d'écrire la relation qui le définit).

On peut maintenant remonter $\dot\tau$ en trois manières, à savoir
\[
\left\{\begin{aligned}
\tau_0&=\sigma_0\widetilde\sigma_0=\widetilde\sigma_0\sigma_0\\
\tau_1&=\sigma_1\widetilde\sigma_1=\widetilde\sigma_1\sigma_1\\
\tau_\infty&=\sigma_\infty\widetilde\sigma_\infty=\widetilde\sigma_\infty\sigma_\infty
\end{aligned}\right.\tag*{(15)}
\]
\[
\text{donc}\quad
\left\{\begin{aligned}
\tau_0^2&=1\\
\tau_1^2&=1\\
\tau_\infty^2&=1
\end{aligned}\right.\tag*{(16)}
\]
